% !TEX root = report.tex
% Настройки ЛР № 2. Подключаются после общей преамбулы ../../shared/preamble.tex
% (неизменённый снимок v2 навыка gost-report); общие настройки здесь не дублируются.

\hypersetup{
    pdftitle={Сравнительный анализ алгоритмов умножения матриц},
    pdfauthor={Дудырев Д.С.}
}

\usepackage{amssymb}
\usepackage{tabularray}
\usepackage{multirow}

% столбцы заданной доли ширины строки: C — по центру, L — по левому краю
\newcolumntype{C}[1]{>{\centering\arraybackslash}m{\dimexpr#1\linewidth-2\tabcolsep\relax}}
\newcolumntype{L}[1]{>{\raggedright\arraybackslash}m{\dimexpr#1\linewidth-2\tabcolsep\relax}}

% числа в таблицах выравниваются по десятичной запятой (столбцы типа S)
\usepackage{siunitx}
\sisetup{output-decimal-marker={,}, group-separator={\,}, group-minimum-digits=4}

\graphicspath{{image/}}

\addto\captionsrussian{\renewcommand{\refname}{СПИСОК ИСПОЛЬЗОВАННЫХ ИСТОЧНИКОВ}}

% заголовки третьего и четвёртого уровней (v2 оформляет только \section и \subsection)
\titleformat{\subsubsection}[hang]
    {\normalfont\normalsize\bfseries\raggedright\hyphenpenalty=10000\exhyphenpenalty=10000}
    {\hspace*{1.25cm}\thesubsubsection}{\fontdimen2\font}{}
\titlespacing*{\subsubsection}{0pt}
    {3.25ex plus 1ex minus .2ex}{1.5ex plus .2ex}
\setcounter{secnumdepth}{4}
\titleformat{\paragraph}[hang]
    {\normalfont\normalsize\bfseries\raggedright\hyphenpenalty=10000\exhyphenpenalty=10000}
    {\hspace*{1.25cm}\theparagraph}{\fontdimen2\font}{}
\titlespacing*{\paragraph}{0pt}
    {3.25ex plus 1ex minus .2ex}{1.5ex plus .2ex}

\newcommand{\floor}[1]{\lfloor #1 \rfloor}

% ---- Строгое соответствие ГОСТ 7.32-2017 (проект algorithm-analysis-2) ----
% п. 6.8.1: выше и ниже каждой формулы не менее одной свободной строки
% (v2 убирает эти отбивки; \normalsize переустанавливает их, поэтому настройка
% добавляется и в \normalsize, и в начало документа)
\makeatletter
\newcommand{\gostdisplayskips}{%
    \setlength{\abovedisplayskip}{\baselineskip}%
    \setlength{\abovedisplayshortskip}{\baselineskip}%
    \setlength{\belowdisplayskip}{\baselineskip}%
    \setlength{\belowdisplayshortskip}{\baselineskip}}
\g@addto@macro\normalsize{\gostdisplayskips}
\AtBeginDocument{\gostdisplayskips}
\makeatother
% п. 6.8.2: пояснение начинается со слова «где» с абзаца
\newcommand{\where}{\par\noindent\hspace*{1.25cm}где}

% п. 6.5.7, 6.6.3: «Рисунок 1 — Наименование», «Таблица 1 — Наименование»;
% п. 6.5.8, 6.6.3: многострочное наименование — через один межстрочный интервал
\DeclareCaptionLabelSeparator{gostdash}{~---~}
\captionsetup{labelsep=gostdash,font=singlespacing}

% п. 6.13: каждая запись содержания — абзац, выровненный влево, номер страницы
% соединяется с наименованием отточием; полужирный только в заголовках (п. 6.1.1).
% Вид — как в ЛР № 1 проекта algorithm-analysis: подразделы и пункты сдвинуты
% под наименование вышестоящей записи, наименования одного уровня начинаются
% ровно, перед разделами и структурными элементами — небольшой интервал
\makeatletter
\renewcommand*\l@section{%
    \addpenalty\@secpenalty\addvspace{0.25em}%
    \@dottedtocline{1}{0pt}{1.5em}}
\renewcommand*\l@subsection{\@dottedtocline{2}{1.5em}{2.3em}}
\renewcommand*\l@subsubsection{\@dottedtocline{3}{3.8em}{3.2em}}
\makeatother
% v2 добавляет к номеру формулы отступы по 1 мм (\tagform@); в ссылках в тексте они не нужны
\renewcommand{\eqref}[1]{\textup{(\ref{#1})}}
% длинные адреса в списке источников переносятся после любого символа
\makeatletter
\g@addto@macro\UrlBreaks{\do\a\do\b\do\c\do\d\do\e\do\f\do\g\do\h\do\i\do\j\do\k\do\l\do\m\do\n\do\o\do\p\do\q\do\r\do\s\do\t\do\u\do\v\do\w\do\x\do\y\do\z\do\0\do\1\do\2\do\3\do\4\do\5\do\6\do\7\do\8\do\9}
\makeatother

\lstdefinelanguage{Swift}{
  morekeywords={func, let, var, if, else, guard, return, while, for, in,
    repeat, switch, case, default, break, continue, struct, class, enum,
    import, nil, true, false, inout, where, private, public, static, self,
    package},
  sensitive=true,
  morecomment=[l]{//},
  morecomment=[s]{/*}{*/},
  morestring=[b]",
}
\lstset{
    inputpath=listings,        % фрагменты кода готовит gen-listings.py
    language=Swift,
    captionpos=b,
    escapeinside={\#*}{*)},    % кириллица набирается самим LaTeX (см. gen-listings.py)
    extendedchars=true,
    abovecaptionskip=4pt,
    % в listings под XeLaTeX пробел после кириллицы теряется без этой подстановки
    % прямая кавычка вместо типографской ” в коде
    literate={\ }{{\ }}1 {"}{{\textquotedbl}}1
}
% v2 запрещает переносы (\hyphenpenalty=10000), а listings переносит длинные
% строки через \discretionary; внутри листингов запрет снимается
\makeatletter
\lst@AddToHook{Init}{\hyphenpenalty=0\exhyphenpenalty=0}
\makeatother
